\chapter*{Abstract}
\paragraph{}This report presents the design and implementation of Hubly, a web application developed as part of the course \textit{Project and Seminar}. The project addresses the lack of centralized information in influencer marketing workflows, where companies and digital content creators often rely on separate social networks, messaging platforms and email channels to discover contacts, negotiate collaborations and track partnership status. This fragmentation makes it difficult to compare opportunities, preserve communication context and manage multiple outreach processes in a consistent way. The proposed solution, Hubly, aims to address this limitation providing a centralized environment where users can register either as companies or as creators. Creator accounts are organized through platform-specific social profiles (for example, a YouTube or Instagram account), allowing the same creator to manage multiple profiles according to their real presence across platforms. Each social channel is represented independently, with its own audience information, content sectors and commercial conditions. Companies and creators can search for each other using criteria such as sector, platform, audience size, price range, company size and country, making the discovery process more structured and aligned with each user’s objectives. In addition, Hubly provides internal conversations associated with the relevant company or creator social profile, ensuring that communication remains linked to the context of a potential partnership. To support organizational needs, Hubly introduces collaborative team management, enabling multiple users from the same organization to manage and interact with the platform under a shared corporate account. Furthermore, to ensure transparency and system accountability, an audit logging mechanism was implemented, providing a record of critical operations performed within the application. The resulting prototype delivers a robust environment for managing company-creator interactions, reducing dependence on external channels and facilitating more organized, scalable, and transparent collaboration workflows.